\documentclass[11pt,a4paper]{article}
\usepackage[utf8]{inputenc}
\usepackage[T1]{fontenc}
\usepackage{amsfonts}
\usepackage[french]{babel}
\frenchbsetup{StandardLists=true} % à inclure si on utilise \usepackage[francais]{babel}
\usepackage{enumitem}
\usepackage{amssymb}
\def\nbR{\ensuremath{\mathrm{I\! R}}}
\usepackage{amsmath}
\DeclareMathOperator{\e}{e}

\usepackage[left=2cm,right=2cm,top=2cm,bottom=2cm]{geometry}
\usepackage{multirow}
\usepackage[dvipsnames]{xcolor}
\usepackage{parskip}
\usepackage{caption}

\usepackage{amsthm}
\newtheorem{theorem}{Théorème}
\newtheorem{corollary}{Corollaire}




\usepackage[T1]{fontenc}
\usepackage{fouriernc}

\usepackage{graphicx}
\usepackage{setspace}
\singlespacing

\usepackage{tcolorbox}
\usepackage{framed}
\usepackage{multicol}

\usepackage{fancybox}
\usepackage{fancyhdr}
\pagestyle{fancy}
\setlength{\headheight}{13.6pt}
\addtolength{\topmargin}{-1.6pt}
\renewcommand\headrulewidth{1pt}
\fancyhead[L]{Lycée Denis Diderot }
\fancyhead[R]{Classe 1BTS}
\setcounter{page}{1}\fancyfoot[C]{Chapitre 4 Cours - Page \thepage}
\fancyfoot[R]{Année 2025 - 2026}
\fancyfoot[L]{Alexis RUIZ}
\usepackage{pstricks,pst-plot}
\usepackage{tikz}
\usepackage{tkz-tab}
\usetikzlibrary{arrows}
\usepackage{pifont}

\newcommand{\ud}{\mathrm{d}}
\newcommand{\V}{\overrightarrow}
\newcommand{\Rep}{(O;\V{u};\V{v})} 
\newcommand{\cadreblanc}[1]{%
\medskip\framebox[\linewidth][c]{%
\rule{0mm}{#1mm}}\par
\medskip}
\newcommand{\cor}[1]{%
\begin{tcolorbox}[colback=red!5!white,colframe=red!75!black]
#1
\end{tcolorbox}
}



\setlength{\columnseprule}{.25mm}
\usepackage{multirow,tabularx,booktabs}
\usepackage{colortbl}
\usepackage{wrapfig}

\usepackage[tikz]{bclogo} 
\usepackage{pgfplots,mathtools}
\usepackage{awesomebox}
\setlength{\aweboxvskip}{0mm}
\newcommand{\defbox}{\awesomebox[black]{2pt}{\faEye}{black}}


\newcommand{\cadre}[1]
{\begin{bclogo}[couleur = white,logo=,  arrondi = 0.3,barre=none]
{%
\rule{0mm}{#1mm}}\par
\medskip
\end{bclogo}}



\begin{document}



\newtcolorbox{mybox}{colback=blue!10!white,
colframe=black!5!black}
\begin{mybox}
\begin{center}    
\textbf{\Large CHAPITRE 4 - DÉRIVATION}\end{center}
\end{mybox}
\vfill
  \begin{bclogo}[couleur = blue!10,logo=\bcinfo,  arrondi = 0.1]
{~Objectifs~:}
~~
\begin{itemize}

\item Calculer la fonction dérivée d'une fonction.
\item Étudier le sens de variation d'une fonction.
\item Exploiter le tableau de variation d'une fonction
\item Mettre en œuvre un procédé pour déterminer la valeur approchée d'une racine d'une équation 
    
    
\end{itemize}
~~
\end{bclogo}

\vfill


\begin{mybox}
\begin{center}    
\textbf{1. Les dérivées usuelles}\end{center}
\end{mybox}
\vfill




\begin{center}
\renewcommand{\arraystretch}{2.6}

\begin{tabularx}{0.7\linewidth}{|c||>{\centering\arraybackslash}X|}
\hline 
Fonctions & Fonctions dérivées \\ 
\hline 
$a$ (constante) & \textcolor{red}{\textbf{0}} \\ 
\hline 
$x^n$ & \textcolor{red}{\textbf{$nx^{n-1}$}} \\ 
\hline 
$\e^x$ & \textcolor{red}{\textbf{$\e^x$}} \\ 
\hline 
$\ln(x)$ & \textcolor{red}{\textbf{$\dfrac{1}{x}$}} \\ 
\hline 
$a^x$ ($a>0$) & \textcolor{red}{\textbf{$a^x\ln(a)$}} \\ 
\hline 
$\sin(x)$ & \textcolor{red}{\textbf{$\cos(x)$}} \\ 
\hline 
$\cos(x)$ & \textcolor{red}{\textbf{$-\sin(x)$}} \\ 
\hline 
$\tan(x)$ & \textcolor{red}{\textbf{$\dfrac{1}{\cos^2 x}$}} \\ 
\hline 
$\arctan(x)$ & \textcolor{red}{\textbf{$\dfrac{1}{1+x^2}$}} \\ 
\hline
\end{tabularx}



\end{center}

\vfill
~~

%%%%%%%%%%%%%%%%%%%%%%%%%
\newpage
%%%%%%%%%%%%%%%%%%%%%%%%%

\fcolorbox{black}{orange!30!}{\textbf{{A1.Dérivation de $x^n$ et $a^x$}}}\\[5pt]
Calculer les fonctions dérivées suivantes : \vfill

\begin{center}
\renewcommand{\arraystretch}{3.9}

\begin{tabularx}{0.7\linewidth}{|c||>{\centering\arraybackslash}X|}
\hline 
Fonctions & Fonctions dérivées \\ 
\hline 
$\dfrac{1}{x}$ & \textcolor{red}{\textbf{$-\dfrac{1}{x^2}$}} \\ 
\hline 
$\dfrac{1}{x^2}$ & \textcolor{red}{\textbf{$-\dfrac{2}{x^3}$}} \\ 
\hline 
$\sqrt{x}$ & \textcolor{red}{\textbf{$\dfrac{1}{2\sqrt{x}}$}} \\ 
\hline 
$3^x$ & \textcolor{red}{\textbf{$3^x\ln 3$}} \\ 
\hline 
$x^{\frac{3}{2}}$ & \textcolor{red}{\textbf{$\dfrac{3}{2}x^{\frac{1}{2}}$}} \\ 
\hline 
$x^4$ & \textcolor{red}{\textbf{$4x^3$}} \\ 
\hline 
$4^x$ & \textcolor{red}{\textbf{$4^x\ln 4$}} \\ 
\hline 
$3x^2-2x+1$ & \textcolor{red}{\textbf{$6x-2$}} \\ 
\hline 
$\dfrac{x^2}{2}+x+1$ & \textcolor{red}{\textbf{$x+1$}} \\ 
\hline 
${x^3}+3x^2+x+1$ & \textcolor{red}{\textbf{$3x^2+6x+1$}} \\ 
\hline 
$x^2+3x$ & \textcolor{red}{\textbf{$2x+3$}} \\ 
\hline 
\end{tabularx}


\end{center}

\vfill
~~

%%%%%%%%%%%%%%%%%%%%%%%%%%%%%%%%%%%%%%%%
\newpage
%%%%%%%%%%%%%%%%%%%%%%%%%%%%%%%%%%%%%%%%

\begin{mybox}
\begin{center}    
\textbf{2. Les formules de dérivation}\end{center}
\end{mybox}
\vfill

\fcolorbox{black}{blue!10!}{\textbf{2.1. Les formules vues en terminale (ou pas)}}\\
\vfill
Si les fonctions $u$ et $v$ sont deux fonctions dérivables sur le même intervalle $I$. \\[5pt]
 On a les résultats suivants : 
\vfill
\begin{center}
\renewcommand{\arraystretch}{3}

\begin{tabularx}{0.5\linewidth}{|c||>{\centering\arraybackslash}X|}
\hline 
Fonction & Fonction dérivée \\ 
\hline 
$k \cdot u(x)$ & \textcolor{red}{\textbf{$k \cdot u'(x)$}} \\ 
\hline 
$u(x) + v(x)$ & \textcolor{red}{\textbf{$u'(x) + v'(x)$}} \\ 
\hline 
$u(x) \times v(x)$ & \textcolor{red}{\textbf{$u'(x)v(x) + u(x)v'(x)$}} \\ 
\hline 
$\dfrac{u(x)}{v(x)}$ & \textcolor{red}{\textbf{$\dfrac{u'(x)v(x) - u(x)v'(x)}{v(x)^2}$}} \\ 
\hline 
\end{tabularx}

\end{center}
\vspace{5mm}
\fcolorbox{black}{green!10!}{\textbf{Exercice 1: Calculer les dérivées des fonctions suivantes }} \\
\begin{multicols}{3}
\begin{center}
$f(x)=3x^4-2x^3+6x+34$

\begin{tcolorbox}[colback=red!5!white,colframe=red!75!black]
{
Dérivons terme à terme : \\[3mm]

dérivée de $(3x^4)=12x^3$ \\[3mm]
dérivée de $(-2x^3)=-6x^2$ \\[3mm]
dérivée de $(6x)=6$ \\[3mm]
dérivée de $(34)=0$ \\[5mm]

$f'(x)=12x^3-6x^2+6$

 \vspace{10mm} 

}

\end{tcolorbox}
\end{center}


\columnbreak

\begin{center}
$g(x)=\dfrac{x-1}{x+2}$

\begin{tcolorbox}[colback=red!5!white,colframe=red!75!black]
{ Avec 
\[
\begin{aligned}
u(x)=x-1 & ~~~~v(x)=x+2. \\
u'(x)=1 & ~~~~v'(x)=1
\end{aligned}
\]

Donc

\[
\begin{aligned}
g'(x) &= \dfrac{1\cdot(x+2)-(x-1)\cdot 1}{(x+2)^2} \\
      &= \dfrac{x+2-(x-1)}{(x+2)^2}
\end{aligned}
\]


\textbf{\textcolor{red}{$g'(x)=\dfrac{3}{(x+2)^2}$}} 
}
\end{tcolorbox}
\end{center}

\columnbreak

\begin{center}
$h(x)=(x^2+1)(x-3)$



\begin{tcolorbox}[colback=red!5!white,colframe=red!75!black]
{ Avec 
\[
\begin{aligned}
u(x) &= x^2+1 & ~~~~ v(x) &= x-3 \\
u'(x) &= 2x & ~~~~ v'(x) &= 1
\end{aligned}
\]
Donc
\[
\begin{aligned}
h'(x) &= u'(x)v(x) + u(x)v'(x) \\
      &= 2x(x-3) + (x^2+1)\cdot 1 \\
      &= 2x^2 - 6x + x^2 + 1 \\
      &= 3x^2 - 6x + 1
\end{aligned}
\]

\vspace{2mm}

\textbf{\textcolor{red}{$h'(x)=3x^2-6x+1$}} 

\vspace{2mm}
}
\end{tcolorbox}

\end{center}

\end{multicols}

%%%%%%%%%%%%%%%%%%%%%%%%%%%%%%
\newpage
%%%%%%%%%%%%%%%%%%%%%%%%%%%%%%

\fcolorbox{black}{orange!30!}{\textbf{A2. Calculer les dérivées des fonctions suivantes :}} \\
\begin{multicols}{2}
\begin{tcolorbox}[colback=red!5!white,colframe=red!75!black]
{
\[
f_1(x) = x^4 - 3x^2 + 2x - 7
\]

\[
\begin{aligned}
\text{dérivée de } x^4 & = 4x^3 \\
\text{dérivée de } -3x^2 & = -6x\\
\text{dérivée de } 2x & = 2\\
\text{dérivée de } -7 & = 0
\end{aligned}
\]

\vspace{4mm}

\begin{center}

    \textbf{\textcolor{red}{$f_1'(x) = 4x^3 - 6x + 2$}}

\end{center}

\vspace{43mm}

}
\end{tcolorbox}

\begin{tcolorbox}[colback=red!5!white,colframe=red!75!black]
{
\[
f_2(x) = 3x \ln(x)
\]

\[
\begin{aligned}
u(x) &= 3x, & v(x) &= \ln(x) \\
u'(x) &= 3, & v'(x) &= \frac{1}{x}
\end{aligned}
\]

\[
\begin{aligned}
    f_2'(x) & = u'(x)v(x) + u(x)v'(x)  \\
            & = 3 \ln(x) + 3x \cdot \frac{1}{x} \\
            &  = 3\ln(x) + 3
\end{aligned}
\]

\begin{center}
    


\textbf{\textcolor{red}{$f_2'(x) = 3\ln(x) + 3$}}

\end{center}

\vspace{24mm}

}
\end{tcolorbox}
\end{multicols}



\begin{multicols}{2}
\begin{tcolorbox}[colback=red!5!white,colframe=red!75!black]
{
\[
f_3(x) = (2x-1)(3-5x)
\]

\vspace{4mm}

\[
\begin{aligned}
u(x) &= 2x-1, & v(x) &= 3-5x \\
u'(x) &= 2, & v'(x) &= -5
\end{aligned}
\]

\vspace{4mm}


\[
\begin{aligned}
    f_3'(x) & = u'(x)v(x) + u(x)v'(x)  \\
            & = 2(3-5x) + (2x-1)(-5) \\
            & = 6-10x -10x+5 \\
            & = -20x + 11
\end{aligned}
\]

\vspace{4mm}


\begin{center}
\textbf{\textcolor{red}{$f_3'(x) = -20x + 11$}}
\end{center}

\vspace{24mm}


}
\end{tcolorbox}

\begin{tcolorbox}[colback=red!5!white,colframe=red!75!black]
{
\[
f_4(x) = (5-x)\e^x
\]

\vspace{4mm}


\[
\begin{aligned}
u(x) &= 5-x, & v(x) &= \e^x \\
u'(x) &= -1, & v'(x) &= \e^x
\end{aligned}
\]

\vspace{4mm}


\[
\begin{aligned}
    f_4'(x) & = u'(x)v(x) + u(x)v'(x)  \\
            & = -1 \cdot \e^x + (5-x)\e^x \\
            & = -\e^x + (5-x)\e^x \\
            & = (4-x)\e^x
\end{aligned}
\]

\vspace{4mm}


\begin{center}
\textbf{\textcolor{red}{$f_4'(x) = (4-x)\e^x$}}
\end{center}

\vspace{24mm}


}
\end{tcolorbox}
\end{multicols}

\newpage
\begin{multicols}{2}
\begin{tcolorbox}[colback=red!5!white,colframe=red!75!black]
{

\vspace{4mm}


\[
f_5(x) = \dfrac{x-2}{2x+1}
\]

\vspace{4mm}

\[
\begin{aligned}
u(x) &= x-2, & v(x) &= 2x+1 \\
u'(x) &= 1, & v'(x) &= 2
\end{aligned}
\]
\vspace{2mm}
\[
\begin{aligned}
    f_5'(x) & = \dfrac{u'(x)v(x) - u(x)v'(x)}{[v(x)]^2}  \\[2mm]
            & = \dfrac{1(2x+1) - (x-2)(2)}{(2x+1)^2} \\[2mm]
            & = \dfrac{2x+1 - 2x+4}{(2x+1)^2} \\[2mm]
            & = \dfrac{5}{(2x+1)^2}
\end{aligned}
\]
\vspace{2mm}
\begin{center}
\textbf{\textcolor{red}{$f_5'(x) = \dfrac{5}{(2x+1)^2}$}}
\end{center}

\vspace{20mm}

}
\end{tcolorbox}

\begin{tcolorbox}[colback=red!5!white,colframe=red!75!black]
{

\vspace{4mm}

\[
f_6(x) = \dfrac{\ln(x)}{2x}
\]

\vspace{4mm}

\[
\begin{aligned}
u(x) &= \ln(x), & v(x) &= 2x \\
u'(x) &= \dfrac{1}{x}, & v'(x) &= 2
\end{aligned}
\]
\vspace{2mm}
\[
\begin{aligned}
    f_6'(x) & = \dfrac{u'(x)v(x) - u(x)v'(x)}{[v(x)]^2}  \\[2mm]
            & = \dfrac{\dfrac{1}{x}(2x) - \ln(x)(2)}{(2x)^2} \\[2mm]
            & = \dfrac{2 - 2\ln(x)}{4x^2} \\[2mm]
            & = \dfrac{1 - \ln(x)}{2x^2}
\end{aligned}
\]
\vspace{2mm}
\begin{center}
\textbf{\textcolor{red}{$f_6'(x) = \dfrac{1 - \ln(x)}{2x^2}$}}
\end{center}

\vspace{13mm}

}
\end{tcolorbox}
\end{multicols}

\vfill 

\begin{multicols}{2}
\begin{tcolorbox}[colback=red!5!white,colframe=red!75!black]
{

\vspace{4mm}

\[
f_7(x) = \sin(x) \times \cos(x)
\]

\vspace{4mm}


\[
\begin{aligned}
u(x) &= \sin(x), & v(x) &= \cos(x) \\
u'(x) &= \cos(x), & v'(x) &= -\sin(x)
\end{aligned}
\]
\vspace{4mm}
\[
\begin{aligned}
    f_7'(x) & = u'(x)v(x) + u(x)v'(x)  \\[2mm]
            & = \cos(x)\cos(x) + \sin(x)(-\sin(x)) \\[2mm]
            & = \cos^2(x) - \sin^2(x)
\end{aligned}
\]
\vspace{6mm}
\begin{center}
\textbf{\textcolor{red}{$f_7'(x) = \cos^2(x) - \sin^2(x)$}}
\end{center}
\vspace{13mm}
}
\end{tcolorbox}

\begin{tcolorbox}[colback=red!5!white,colframe=red!75!black]
{

\vspace{4mm}

\[
f_8(x) = \cos^2(x)
\]

\vspace{4mm}


\[
\begin{aligned}
u(x) &= \cos(x) \\
u'(x) &= -\sin(x)
\end{aligned}
\]
\vspace{4mm}
\[
\begin{aligned}
    f_8'(x) & = 2u(x) \times u'(x)  \\[2mm]
            & = 2\cos(x) \times (-\sin(x)) \\[2mm]
            & = -2\sin(x)\cos(x)
\end{aligned}
\]
\vspace{6mm}
\begin{center}
\textbf{\textcolor{red}{$f_8'(x) = -2\sin(x)\cos(x)$}}
\end{center}
\vspace{13mm}
}
\end{tcolorbox}
\end{multicols}

\vfill 


%%%%%%%%%%%%%%%%%%%%%%%%%%%%%%
\newpage
%%%%%%%%%%%%%%%%%%%%%%%%%%%%%%


\fcolorbox{black}{green!10!}{\textbf {Exercice 2 }}\\[5 pt]
Déterminer les dérivées des fonctions suivantes

\begin{multicols}{2}
\begin{tcolorbox}[colback=red!5!white,colframe=red!75!black]
{
\textbf{ Fonction $f(x)=\e^x(\e^x+3)$}

\begin{align*}
u(x) & =\e^x  & ~~~~ u'(x) & =\e^x\\
v(x) & =\e^x +3  & ~~~~ v'(x) &  =\e^x
\end{align*}

\begin{align*}
f'(x) &= \e^x \times (\e^x+3) + \e^x \times \e^x \\
f'(x) &= \e^x(2\e^x + 3)
\end{align*}

}
\end{tcolorbox}

\columnbreak

\begin{tcolorbox}[colback=red!5!white,colframe=red!75!black]
{
\textbf{Fonction $g(x)=(\e^x-1)^2=(\e^x-1) \times (\e^x-1)$}

\begin{align*}
u(x) & =\e^x - 1  & ~~~~ u'(x) & =\e^x \\
v(x) & =\e^x - 1  & ~~~~ v'(x) &  =\e^x
\end{align*}

\begin{align*}
g'(x) &= \e^x(\e^x-1) + (e^x -1)\e^x \\
g'(x) &= 2\e^x(\e^x - 1)
\end{align*}

}
\end{tcolorbox}

\end{multicols}


\noindent
\begin{tcolorbox}[colback=red!5!white,colframe=red!75!black]
{
\vspace{1cm}
Si les fonctions $g$ et $u$ sont des fonctions dérivables sur les intervalles considérés on   le résultats important suivant : 
$$(g(u(x)))' = u'(x) \times g'(u(x))$$
\vspace{1cm}
}
\end{tcolorbox}
\noindent

Ce résultat permet d'obtenir les dérivées des fonctions du tableau 1 en ayant remplacé $x$ par $u(x)$. \\[0.2 cm]
On obtient ainsi les dérivées suivantes : \\

\begin{center}
{
\renewcommand{\arraystretch}{2.7}
\begin{tabularx}{1\linewidth}{|>{\centering\arraybackslash}c||>{\centering\arraybackslash}p{3.5cm}|||>{\centering\arraybackslash}X|}
\hline
Fonctions & Fonctions dérivées & Exemples \\
\hline
$\left(\e^u\right)$ & \textcolor{red}{$u'\e^u$} & \textcolor{red}{$\left(\e^{2x}\right)' = 2\e^{2x}$} \\
\hline
$u^n$ & \textcolor{red}{$nu'u^{n-1}$} & \textcolor{red}{$\left((3x+1)^5\right)' = 5 \times 3 \times (3x+1)^4 = 15(3x+1)^4$} \\
\hline
$\ln u$ & \textcolor{red}{$\dfrac{u'}{u}$} & \textcolor{red}{$\left(\ln(2x+3)\right)' = \dfrac{2}{2x+3}$} \\
\hline
$\sin u$ & \textcolor{red}{$u'\cos u$} & \textcolor{red}{$\left(\sin(3x)\right)' = 3\cos(3x)$} \\
\hline
$\cos u$ & \textcolor{red}{$-u'\sin u$} & \textcolor{red}{$\left(\cos(2x)\right)' = -2\sin(2x)$} \\
\hline
$\arctan u$ & \textcolor{red}{$\dfrac{u'}{1+u^2}$} & \textcolor{red}{$\left(\arctan(x^2)\right)' = \dfrac{2x}{1+x^4}$} \\
\hline
\end{tabularx}
}
\end{center}

\newpage

\fcolorbox{black}{green!10!}{\textbf {Exercice 3 }}\\[5 pt]
Déterminer les dérivées des fonctions suivantes
\begin{multicols}{2}
$f(x)=\e^{3x}(2x+1) $
\vfill
\begin{tcolorbox}[colback=red!5!white,colframe=red!75!black]
\begin{itemize}
\item $u(x) = e^{3x}$ donc $u'(x) = 3e^{3x}$
\item $v(x) = 2x+1$ donc $v'(x) = 2$
\end{itemize}

\begin{align*}
f'(x) &= 3e^{3x}(2x+1) + e^{3x} \times 2 \\
&= e^{3x}[6x + 3 + 2] \\
&= e^{3x}(6x + 5)
\end{align*}


$$\boxed{f'(x) = e^{3x}(6x + 5)}$$

\end{tcolorbox}

$g(x)=(x^2-x+1)\e^{-x}$
\vfill

\begin{tcolorbox}[colback=red!5!white,colframe=red!75!black]
\begin{align*}
u(x)& = x^2-x+1 & ~~ &  u'(x) = 2x-1 \\
v(x) & = e^{-x} &  ~~ &  v'(x) = -e^{-x}
\end{align*}

\begin{align*}
g'(x) &= (2x-1)e^{-x} - (x^2-x+1)e^{-x} \\
&= e^{-x}[2x - 1 - x^2 + x - 1] \\
&= e^{-x}(-x^2 + 3x - 2)
\end{align*}

\end{tcolorbox}


\end{multicols}


\fcolorbox{black}{green!10!}{\textbf {Exercice 4 }}\\[5 pt]
Déterminer les dérivées des fonctions suivantes
\begin{multicols}{2}
$f(x)=\dfrac{1}{4}x^2-\dfrac{1}{4}-\dfrac{1}{2}\ln(x)$
\vfill
\cor{\begin{align*}
f'(x) &= \frac{1}{4} \times 2x - \frac{1}{2} \times \frac{1}{x} \\
&= \frac{1}{2}x - \frac{1}{2x} \\
&= \frac{x^2 - 1}{2x}
\end{align*}
}

\columnbreak

$g(x)= \ln\left(1+\e^x \right)$
\cor{$[\ln(u)]' = \dfrac{u'}{u}$ avec $u(x) = 1+e^x$ donc $u'(x) = e^x$.

\begin{align*}
g'(x) &= \frac{e^x}{1+e^x}
\end{align*}

\vspace{0.5cm}
}
\end{multicols}



\fcolorbox{black}{green!10!}{\textbf {Exercice 5 }}\\[5 pt]
Déterminer les dérivées des fonctions suivantes
\begin{multicols}{2}
$f(x)=\dfrac{\e^x-\e^{-x}}{\e^x+\e^{-x}}$
\vfill
\cor{
$\left(\dfrac{u}{v}\right)' = \dfrac{u'v - uv'}{v^2}$ avec :
\begin{align*}
u(x) = & e^x - e^{-x} & u'(x) = & e^x + e^{-x} \\
v(x) = & e^x + e^{-x} & v'(x) = & e^x - e^{-x} \\
f'(x) = & \frac{(e^x + e^{-x})^2 - (e^x - e^{-x})^2}{(e^x + e^{-x})^2}
\end{align*}

Développons :
\begin{itemize}
\item $(e^x + e^{-x})^2 = e^{2x} + 2 + e^{-2x}$
\item $(e^x - e^{-x})^2 = e^{2x} - 2 + e^{-2x}$
\end{itemize}

$$f'(x) = \frac{4}{(e^x + e^{-x})^2}$$
}

$g(x)= \ln\left(\dfrac{1+x}{1-x}\right)$
\cor{
On utilise les propriétés du logarithme :
$$g(x) = \ln(1+x) - \ln(1-x)$$

Donc :
\begin{align*}
g'(x) &= \frac{1}{1+x} - \frac{-1}{1-x} \\
&= \frac{1}{1+x} + \frac{1}{1-x} \\
&= \frac{1-x + 1+x}{(1+x)(1-x)} \\
&= \frac{2}{1-x^2}
\end{align*}
}

\end{multicols}

\newpage

\fcolorbox{black}{green!10!}{\textbf {Exercice 6 }}\\[5 pt]
Déterminer les dérivées des fonctions suivantes
\begin{multicols}{2}
$f(x)=x^2\ln(x)$
\cor{
$(uv)' = u'v + uv'$ avec :
\begin{itemize}
\item $u(x) = x^2$ donc $u'(x) = 2x$
\item $v(x) = \ln(x)$ donc $v'(x) = \dfrac{1}{x}$
\end{itemize}

\begin{align*}
f'(x) &= 2x\ln(x) + x^2 \cdot \frac{1}{x} \\
&= 2x\ln(x) + x \\
&= x(2\ln(x) + 1)
\end{align*}

}
$g(x)= (x^2-1)\e^{2x}$
\cor{
$(uv)' = u'v + uv'$ avec :
\begin{itemize}
\item $u(x) = x^2-1$ donc $u'(x) = 2x$
\item $v(x) = e^{2x}$ donc $v'(x) = 2e^{2x}$
\end{itemize}

\begin{align*}
g'(x) &= 2xe^{2x} + 2(x^2-1)e^{2x} \\
&= e^{2x}[2x + 2x^2 - 2] \\
&= 2e^{2x}(x^2 + x - 1)
\end{align*}
}


\end{multicols}

\fcolorbox{black}{green!10!}{\textbf {Exercice 7 }}\\[5 pt]
Déterminer les dérivées des fonctions suivantes
\begin{multicols}{2}

$f(x)=\e^{\frac{2x-1}{2-x}}$
\cor{
$[e^u]' = u'e^u$ avec $u(x) = \dfrac{2x-1}{2-x}$.

Calculons $u'(x)$ :
\begin{align*}
u'(x) &= \frac{2(2-x) - (2x-1)(-1)}{(2-x)^2} \\
&= \frac{4 - 2x + 2x - 1}{(2-x)^2} \\
&= \frac{3}{(2-x)^2}
\end{align*}

Donc :
$$f'(x) = \frac{3}{(2-x)^2} \cdot e^{\frac{2x-1}{2-x}}$$

$$\boxed{f'(x) = \frac{3e^{\frac{2x-1}{2-x}}}{(2-x)^2}}$$
}

\columnbreak 


$g(x)= \ln\left(\dfrac{e^x}{e^x+1}\right)$
\cor{
On utilise les propriétés du logarithme :
$$g(x) = \ln(e^x) - \ln(e^x+1) = x - \ln(e^x+1)$$

Donc :
\begin{align*}
g'(x) &= 1 - \frac{e^x}{e^x+1} \\
&= \frac{e^x+1 - e^x}{e^x+1} \\
&= \frac{1}{e^x+1}
\end{align*}

$$\boxed{g'(x) = \frac{1}{e^x+1}}$$
}


\end{multicols}

%%%%%%%%%%%%%%%%%%%%%%%%%%%%%%%%%%%
\newpage           %%%%%%%%%%%%%%%%
%%%%%%%%%%%%%%%%%%%%%%%%%%%%%%%%%%%

\fcolorbox{black}{blue!10!}{\textbf {2.2 Dérivées successives  }}\\[5 pt]
\cor{
Soit $f$ une fonction dérivable sur un intervalle $I$.

\begin{itemize}
\item La \textbf{dérivée première} de $f$ est notée $f'$ ou $f^{(1)}$.

\item La \textbf{dérivée seconde} de $f$ est la dérivée de $f'$, notée $f''$ ou $f^{(2)}$.

\item La \textbf{dérivée troisième} de $f$ est la dérivée de $f''$, notée $f'''$ ou $f^{(3)}$.

\item Plus généralement, la \textbf{dérivée $n$-ième} de $f$ est la dérivée de $f^{(n-1)}$, notée $f^{(n)}$ pour $n \in \mathbb{N}^*$.
\end{itemize}

\vfill 

On dit que $f$ est de \textbf{classe $\mathcal{C}^n$} sur $I$ si $f$ admet des dérivées jusqu'à l'ordre $n$ et si $f^{(n)}$ est continue sur $I$.

\vfill 

\textbf{Notations alternatives :} ~~~~~~$f^{(n)}(x) = \frac{d^n f}{dx^n}(x) = D^n f(x)$
}

\vfill

\cor{
\begin{multicols}{2}

    avec $f(x)=x^4$ 
    
\begin{align*}
    f^{(1)}(x) & =4x^3 \\
    f^{(2)}(x) & =12x^2 
\end{align*}
\begin{align*}
    f^{(3)}(x) & =24x \\
    f^{(4)}(x) & =24 \\
    f^{(5)}(x) & =0 \\
\end{align*}
\end{multicols}


}

\fcolorbox{black}{orange!30!}{\textbf {A3. Déterminer $f^{(5)}$ avec $f(x)=\cos(3x)$}}\\[5 pt]
\cor{
\begin{align*}
f(x) &= \cos(3x) & ~~\\
f'(x) &= -3\sin(3x) &  ~~\\
f''(x) &= -3^2\cos(3x) & = -9\cos(3x) \\
f^{(3)}(x) &= 3^3\sin(3x) & = 27\sin(3x) \\
f^{(4)}(x) &= 3^4\cos(3x) & = 81\cos(3x) \\
f^{(5)}(x) &= -3^5\sin(3x) & = -243\sin(3x)
\end{align*}
}

\vfill 

\fcolorbox{black}{blue!10!}{\textbf {3. Les principales utilisations de la dérivée  }}\\[5 pt]

\vfill 

\fcolorbox{black}{blue!10!}{\textbf {3.1 Équation des tangentes et variations d'une fonction  }}\\[5 pt]

\cor{
En classe de Terminale vous avez vu l'équation de la tangente au point d'abscisse $a$ , à la courbe représentant la fonction $f$. 

\vspace{0.5cm}
\fbox{$y=f'(a)(x-a)+f(a)$}

En développant cette expression on retrouve le résultat vu en première : la tangente à la courbe représentant la fonction $f$ au point d'abscisse $a$ a pour coefficient directeur $f'(a)$. 

\vspace{0.5cm}

}

\vfill 

%%%%%%%%%%%%%%%%%%%%%%%%%%%%%%%
\newpage
%%%%%%%%%%%%%%%%%%%%%%%%%%%%%%%

\fcolorbox{black}{orange!30!}{\textbf {A4. Une équation de la tangente}}\\[5 pt]
Soit $f$ une fonction définie sur $\mathbb{R}$ par $f(x)=(2x+1)\e^x$ et C sa courbe représentative. 
Donner une équation de la tangente à C au point d'abscisse $x=0$.Vérifier votre résultat à l'aide de la calculatrice en traçant la courbe représentative de $f$ et la tangente trouvée. \\[5 pt]
\cor{
\begin{multicols}{2}
  L'équation de la tangente à $\mathcal{C}$ au point d'abscisse $x_0$ est donnée par :$y = f(x_0) + f'(x_0)(x - x_0)$

Avec, $x_0 = 0$, $y = f(0) + f'(0) \cdot x$

 Calcul de $f(0)$
\begin{align*}
f(0) &= (2 \times 0 + 1)e^0 & = 1 \times 1  & = 1
\end{align*}

Calcul de $f'(x)$

On utilise la formule du produit $(uv)' = u'v + uv'$ avec :
\begin{itemize}
\item $u(x) = 2x+1$ donc $u'(x) = 2$
\item $v(x) = e^x$ donc $v'(x) = e^x$
\end{itemize}

\begin{align*}
f'(x) &= 2 \cdot e^x + (2x+1) \cdot e^x \\
&= e^x[2 + 2x + 1] \\
&= e^x(2x + 3)
\end{align*}

\textbf{Étape 3 :} Calcul de $f'(0)$

\begin{align*}
f'(0) &= e^0(2 \times 0 + 3) \\
&= 1 \times 3 \\
&= 3
\end{align*}

\textbf{Étape 4 :} Équation de la tangente

\begin{align*}
y &= f(0) + f'(0) \cdot x \\
y &= 1 + 3x
\end{align*}

\cor{
L'équation de la tangente à $\mathcal{C}$ au point d'abscisse $x=0$ est :
$$\boxed{y = 3x + 1}$$

Le point de tangence est $A(0, 1)$ et le coefficient directeur de la tangente est $m = 3$.
}
  
\end{multicols}
}

\vfill

\fcolorbox{black}{blue!10!}{\textbf {3.2 Théorèmes relatifs au sens de variation d'une fonction }}\\[5 pt]
\defbox{\fbox{
\begin{minipage}{14cm}{

Si sur un intervalle $I$ la dérivée $f'$ de la fonction $f$ est positive alors la fonction $f$ est croissante sur $I$.\\[5 pt]
Si sur un intervalle $I$ la dérivée $f'$ de la fonction $f$ est négative alors la fonction $f$ est décroissante sur $I$.
}

\end{minipage}
}}

\vfill

\fcolorbox{black}{green!10!}{\textbf {Exercice 8 - Faire le point avec un QCM}}\\[5 pt]
Soit $f$ la fonction définie sur $\mathbb{R}$ par $f(x)=\e^{2x}(x^2+x-1)$ et $C$ sa courbe représentative : \\
\begin{multicols}{2}

1. Dérivée 
\begin{itemize}[label=$\square$]
\item   a) $f'(x)=(x^2+3x)\e^{2x}$
\item   \textcolor{red}{$\boxtimes$ b) $f'(x)=(2x^2+4x-1)\e^{2x}$}
\item   c) $f'(x)=(2x+1)\e^{2x}$
\end{itemize}
\begin{center}
\noindent\rule{5cm}{0.4pt}
\end{center}

\textbf{2. Une équation de la droite tangente à $C$ au point d'abscisse $a=0$ est :}
\begin{itemize}[label=$\square$]
\item   a) $y=3x+1$
\item   b) $y=-1$
\item   \textcolor{red}{$\boxtimes$ c) $y=-x-1$} 
\end{itemize}


\textbf{3. $f$ est strictement croissante sur :}
\begin{itemize}[label=$\square$]
\item   a) $\mathbb{R}$
\item   \textcolor{red}{$\boxtimes$ b) $[1;+\infty[$}
\item   c) $[0;1]$
\end{itemize}

\begin{center}
\noindent\rule{5cm}{0.4pt}
\end{center}

\textbf{4. $f$ est strictement positive sur :}  
\begin{itemize}[label=$\square$]
\item   a) $\mathbb{R}$
\item   \textcolor{red}{$\boxtimes$ b) $[1;+\infty[$}
\item   c) $[0;1]$
\end{itemize}

\end{multicols}

\vfill 
\newpage 
%%%%%%%%%%%%%%%%%%%%%%%%%%%%%%%%%%%%%%%


\fcolorbox{black}{green!10!}{ \textbf {Exercice 9 - Avec plusieurs fonctions dérivées}}\\
\normalsize
Soit la fonction $f$ définie sur $\mathbf{R}$ par $f(x)=(3x-1)\e^{2x}$ : \\

\textbf{1. Calculer $f'(x)$ et $f''(x).$}

\begin{multicols}{2}
\begin{tcolorbox}[colback=red!5!white,colframe=red!75!black]
{
\textbf{Calcul de $f'(x)$ :}
\[
f(x) = (3x-1)\e^{2x}
\]
\[
\begin{aligned}
u(x) &= 3x-1 & u'(x) &= 3 \\
v(x) &= \e^{2x} & v'(x) &= 2\e^{2x}\\
\end{aligned}
\]
\[
\begin{aligned}
    f'(x) & = u'(x)v(x) + u(x)v'(x)  \\
            & = 3\e^{2x} + (3x-1) \cdot 2\e^{2x} \\
            & = \e^{2x}[3 + 6x - 2] \\
\end{aligned}
\]
\begin{center}
\textbf{\textcolor{red}{$f'(x) = \e^{2x}(6x + 1)$}}
\end{center}
}
\end{tcolorbox}

\begin{tcolorbox}[colback=red!5!white,colframe=red!75!black]
{
\textbf{Calcul de $f''(x)$ :}
\[
f'(x) = (6x+1)\e^{2x}
\]
\[
\begin{aligned}
u(x) &= 6x+1 & u'(x) &= 6 \\
v(x) &= \e^{2x} & v'(x) &= 2\e^{2x}
\end{aligned}
\]
\[
\begin{aligned}
    f''(x) & = u'(x)v(x) + u(x)v'(x)  \\
            & = 6\e^{2x} + (6x+1) \cdot 2\e^{2x} \\
            & = \e^{2x}[6 + 12x + 2] 
\end{aligned}
\]
\begin{center}
\textbf{\textcolor{red}{$f''(x) = \e^{2x}(12x + 8)$}}
\end{center}
}
\end{tcolorbox}
\end{multicols}

\textbf{2. Montrer que pour tout réel $x$, $f''(x) - f'(x) - 2f(x) = 9\e^{2x}.$}

\begin{tcolorbox}[colback=red!5!white,colframe=red!75!black]
{
\[
\begin{aligned}
f''(x) - f'(x) - 2f(x) &= \e^{2x}(12x+8) - \e^{2x}(6x+1) - 2(3x-1)\e^{2x} \\
&= \e^{2x}[(12x+8) - (6x+1) - 2(3x-1)] \\
&= \e^{2x}[12x + 8 - 6x - 1 - 6x + 2] \\
&= \e^{2x}[(12x - 6x - 6x) + (8 - 1 + 2)] \\
&= \e^{2x}[0 + 9] \\
&= 9\e^{2x}
\end{aligned}
\]

\vspace{2mm}

\begin{center}
\textbf{\textcolor{red}{On a bien démontré que $f''(x) - f'(x) - 2f(x) = 9\e^{2x}$}}
\end{center}
}
\end{tcolorbox}

\fcolorbox{black}{green!10!}{\textbf {Exercice 10 - Équation de la tangente}}\\

Écrire une équation de la tangente à la courbe représentative de la fonction $f$ au point d'abscisse $a$ donné : 

\begin{multicols}{2}
\begin{tcolorbox}[colback=red!5!white,colframe=red!75!black]
{
\textbf{$f(x)= \ln x + 2x - 3$ en $a=1$}

\vspace{3mm}

\textbf{Calcul de $f(1)$ :}
$
f(1) = \ln(1) + 2(1) - 3 = -1
$

\vspace{2mm}

\textbf{Calcul de $f'(x)$ :}
$f'(x) = \dfrac{1}{x} + 2$


\vspace{2mm}

\textbf{Calcul de $f'(1)$ :}
$
f'(1) = \dfrac{1}{1} + 2 = 1 + 2 = 3
$

\vspace{2mm}

\textbf{Équation de la tangente :}
\[
\begin{aligned}
y &= f'(1)(x-1) + f(1) \\
y &= 3(x-1) + (-1) \\
y &= 3x - 3 - 1 \\
\end{aligned}
\]
\begin{center}
\textbf{\textcolor{red}{$y = 3x - 4$}}
\end{center}
}
\end{tcolorbox}

\begin{tcolorbox}[colback=red!5!white,colframe=red!75!black]
{
\textbf{$f(x)=\e^x(3-2x)$ en $a=0$}

\textbf{Calcul de $f(0)$ :}
$f(0) = \e^0(3-0) = 1 \times 3 = 3$


\vspace{2mm}

\textbf{Calcul de $f'(x)$ :}
\[
\begin{aligned}
f'(x) &= \e^x(3-2x) + \e^x(-2) \\
&= \e^x(3-2x-2) \\
&= \e^x(1-2x)
\end{aligned}
\]

\vspace{2mm}

\textbf{Calcul de $f'(0)$ :}
$f'(0) = \e^0(1-0) = 1 \times 1 = 1$


\vspace{2mm}

\textbf{Équation de la tangente :}
\[
\begin{aligned}
y &= f'(0)(x-0) + f(0) \\
y &= 1 \times x + 3 \\
\end{aligned}
\]
\begin{center}
\textbf{\textcolor{red}{$y = x + 3$}}
\end{center}
}
\end{tcolorbox}
\end{multicols}

%%%%%%%%%%%%%%%%%%%%%%%%%%%%%%%%%%%%%%
\newpage
%%%%%%%%%%%%%%%%%%%%%%%%%%%%%%%%%%%%%%

    
 \fcolorbox{black}{blue!10!}{\textbf {3.3. Le théorème des valeurs intermédiaires.   }}

 \vfill 

 \cor{

Soit $f$ une fonction continue sur un intervalle $[a,b]$ avec $a < b$.

Si $f(a) \leq k \leq f(b)$ (ou $f(b) \leq k \leq f(a)$), alors il existe au moins un réel $c \in [a,b]$ tel que $f(c) = k$.
}

\vfill 

 \begin{tcolorbox}[colback=red!5!white,colframe=red!75!black]
Soit $f$ une fonction continue sur un intervalle $[a,b]$.

Si $f(a)$ et $f(b)$ sont de signes opposés, alors l'équation $f(x) = 0$ admet au moins une solution dans $]a,b[$.
\end{tcolorbox}

\vfill 

\textbf {Exemple  }\\[0.1cm]

Étudier le sens de variation de la fonction $f$ définie sur $\mathbf{R}$ par $f(x) = \e^{-x}(x+2)$ et en déduire sans 
calcul que l'équation $f(x)=0$ n'admet qu'une solution dans $\mathbf{R}$.


\begin{tcolorbox}[colback=red!5!white,colframe=red!75!black]
\begin{multicols}{2}
La fonction $f$ est définie sur $\mathbb{R}$ par :
$$f(x) = e^{-x}(x+2)$$

En utilisant la formule de dérivation d'un produit $(uv)' = u'v + uv'$ :

\begin{align*}
f'(x) &= \left(e^{-x}\right)'(x+2) + e^{-x}(x+2)' \\
f'(x) &= -e^{-x}(x+2) + e^{-x} \times 1 \\
f'(x) &= e^{-x}\left[-(x+2) + 1\right] \\
f'(x) &= e^{-x}(-x-1) \\
f'(x) &= -e^{-x}(x+1)
\end{align*}

\subsection*{3) Calcul des limites(calculatrice)}

$$\lim_{x \to -\infty} f(x) =-\infty ~~\lim_{x \to +\infty} f(x) =  0^+$$

\subsection*{4) Étude du signe de $f'(x)$}

On a $f'(x) = -e^{-x}(x+1)$.

Puisque $e^{-x} > 0$ pour tout $x \in \mathbb{R}$, on a :

\begin{align*}
f'(x) > 0 &\iff -e^{-x}(x+1) > 0 \\
&\iff -(x+1) > 0 \\
&\iff x+1 < 0 \\
&\iff x < -1
\end{align*}

\subsection*{3) Valeur au maximum}

$$f(-1) = e^{-(-1)}(-1+2) = e^1 \times 1 = e$$

\subsection*{5) Tableau de variation}

\begin{center}
\begin{tabular}{|c|ccccc|}
\hline
$x$ & $-\infty$ & & $-1$ & & $+\infty$ \\
\hline
$f'(x)$ & & $+$ & $0$ & $-$ & \\
\hline
& & & $e$ & & \\
$f(x)$ & $-\infty$ & $\nearrow$ & & $\searrow$ & $0^+$ \\
& & & & & \\
\hline
\end{tabular}
\end{center}

\vspace{0.5cm}

La fonction $f$ est \textbf{strictement croissante} sur $]-\infty, -1]$ et \textbf{strictement décroissante} sur $[-1, +\infty[$ avec un maximum en $x = -1$ où $f(-1) = e > 0$.

\subsection*{6) Conclusion sur l'équation $f(x) = 0$}

\begin{itemize}
\item Sur $]-\infty, -1]$ : la fonction $f$ part de $-\infty$ et atteint $e > 0$. Par continuité et stricte croissance, elle traverse l'axe des abscisses \textbf{exactement une fois}. Il existe donc une unique valeur $\alpha \in ]-\infty, -1[$ telle que $f(\alpha) = 0$.

\item Sur $[-1, +\infty[$ : la fonction $f$ part de $e > 0$ et tend vers $0^+$ en restant strictement positive. Elle ne s'annule donc \textbf{jamais} sur cet intervalle.
\end{itemize}


\textbf{Conclusion :} L'équation $f(x) = 0$ admet une unique solution dans $\mathbb{R}$, située dans l'intervalle $]-\infty, -1[$.
\end{multicols}
\end{tcolorbox}

\vfill 

%%%%%%%%%%%%%%%%%%%%%%%%%%%%%%%%%%%%%%%%%%
\newpage
%%%%%%%%%%%%%%%%%%%%%%%%%%%%%%%%%%%%%%%%%%


\fcolorbox{black}{green!10!}{\textbf {Exercice 11 :} }\\ [5 pt]
Soit la fonction $f$ définie sur $\mathbf{R}$ par $f(x)=\e^x(4-2x)$

On nomme C la courbe représentative dans un repère orthonormé.

\textbf{1. Déterminer les limites de $f$ en $-\infty$ et  $+\infty$.}

\begin{tcolorbox}[colback=red!5!white,colframe=red!75!black]
\begin{multicols}{2}
{

\begin{center}

\textbf{Limite en $-\infty$ :}


\textbf{\textcolor{red}{$\displaystyle \lim_{x \to -\infty} f(x) = 0$}}

\columnbreak

\textbf{Limite en $+\infty$ :}


\textbf{\textcolor{red}{$\displaystyle \lim_{x \to +\infty} f(x) = -\infty$}}

\end{center}
}
\end{multicols}
\end{tcolorbox}

\vspace{3mm}

\textbf{2. Calculer $f'(x)$ et étudier son signe sur $\mathbf{R}$}

\begin{tcolorbox}[colback=red!5!white,colframe=red!75!black]
{
\begin{multicols}{2}
\textbf{Calcul de $f'(x)$ :}
\[
f(x) = \e^x(4-2x)
\]
\[
\begin{aligned}
u(x) &= \e^x, & v(x) &= 4-2x \\
u'(x) &= \e^x, & v'(x) &= -2
\end{aligned}
\]
\[
\begin{aligned}
    f'(x) & = u'(x)v(x) + u(x)v'(x)  \\
            & = \e^x(4-2x) + \e^x(-2) \\
            & = \e^x[4-2x-2] \\
            & = \e^x(2-2x)
\end{aligned}
\]
\begin{center}
\textbf{\textcolor{red}{$f'(x) = \e^x(2-2x)$}}
\end{center}
\end{multicols}
}
\end{tcolorbox}

\begin{tcolorbox}[colback=red!5!white,colframe=red!75!black]
\begin{multicols}{2}
{
\textbf{Étude du signe de $f'(x)$ :}

$\e^x > 0$ pour tout $x \in \mathbf{R}$

Le signe de $f'(x)$ dépend donc de $(2-2x)$ :
\[
2-2x = 0 \Longleftrightarrow x = 1
\]

\vspace{2mm}

\begin{itemize}
    \item Si $x < 1$ : $2-2x > 0$ donc \textbf{\textcolor{red}{$f'(x) > 0$}}
    \item Si $x = 1$ : \textbf{\textcolor{red}{$f'(x) = 0$}}
    \item Si $x > 1$ : $2-2x < 0$ donc \textbf{\textcolor{red}{$f'(x) < 0$}}
\end{itemize}
}
\end{multicols}
\end{tcolorbox}

\vspace{3mm}

\textbf{3. Donner le tableau de variation de $f$.}

\begin{tcolorbox}[colback=red!5!white,colframe=red!75!black]
{
\[
f(1) = \e^1(4-2 \times 1) = \e(4-2) = 2\e
\]

\begin{center}
\begin{tikzpicture}
    \tkzTabInit[lgt=2,espcl=2.5]
    {$x$ /1, $f'(x)$ /1, $f(x)$ /2}
    {$-\infty$, $1$, $+\infty$}
    \tkzTabLine{, +, z, -, }
    \tkzTabVar{-/$0$, +/$2\e$, -/$-\infty$}
\end{tikzpicture}
\end{center}

}
\end{tcolorbox}

\vspace{3mm}

\textbf{4. Écrire une équation de la tangente à C au point d'abscisse $a=1$}

\begin{tcolorbox}[colback=red!5!white,colframe=red!75!black]
\begin{multicols}{2}
{
L'équation de la tangente au point d'abscisse $a$ est : $y = f'(a)(x-a) + f(a)$

\vspace{2mm}

Pour $a = 1$ :
\[
\begin{aligned}
f(1) &= \e^1(4-2) = 2\e \\
f'(1) &= \e^1(2-2) = 0
\end{aligned}
\]

\vspace{2mm}

\[
\begin{aligned}
y &= f'(1)(x-1) + f(1) \\
y &= 0(x-1) + 2\e \\
y &= 2\e
\end{aligned}
\]

\vspace{2mm}

\begin{center}
\textbf{\textcolor{red}{L'équation de la tangente est : $y = 2\e$}}
\end{center}

(Tangente horizontale)
}
\end{multicols}
\end{tcolorbox}

%%%%%%%%%%%%%%%%%%%%%%%%%%%%%%%%%%
\newpage
%%%%%%%%%%%%%%%%%%%%%%%%%%%%%%%%%%

\fcolorbox{black}{green!10!}{\textbf {Exercice 12 :} }\\ [5 pt]
Soit la fonction $f$ définie sur $\mathbf{R}$ par $f(x)=\e^{-x}(3-4x)$

On nomme C la courbe représentative dans un repère orthonormé.

\textbf{1. Déterminer les limites de $f$ en $-\infty$ et  $+\infty$.}

\begin{tcolorbox}[colback=red!5!white,colframe=red!75!black]
\begin{multicols}{2}
    

{
\begin{center}
    

\textbf{Limite en $-\infty$ :}

 \textbf{\textcolor{red}{$\displaystyle \lim_{x \to -\infty} f(x) = +\infty$}}
\columnbreak

\textbf{Limite en $+\infty$ :}

 \textbf{\textcolor{red}{$\displaystyle \lim_{x \to +\infty} f(x) = 0$}}
\end{center}}
\end{multicols}
\end{tcolorbox}

\vspace{3mm}

\textbf{2. Calculer $f'(x)$ et étudier son signe sur $\mathbf{R}$}

\begin{tcolorbox}[colback=red!5!white,colframe=red!75!black]
{

\begin{multicols}{2}
    

\textbf{Calcul de $f'(x)$ :}
\[
f(x) = \e^{-x}(3-4x)
\]
\[
\begin{aligned}
u(x) &= \e^{-x}, & v(x) &= 3-4x \\
u'(x) &= -\e^{-x}, & v'(x) &= -4
\end{aligned}
\]
\[
\begin{aligned}
    f'(x) & = u'(x)v(x) + u(x)v'(x)  \\
            & = -\e^{-x}(3-4x) + \e^{-x}(-4) \\
            & = \e^{-x}[-3+4x-4] \\
            & = \e^{-x}(4x-7)
\end{aligned}
\]
\begin{center}
\textbf{\textcolor{red}{$f'(x) = \e^{-x}(4x-7)$}}
\end{center}

\end{multicols}

}
\end{tcolorbox}

\begin{tcolorbox}[colback=red!5!white,colframe=red!75!black]

\begin{multicols}{2}
    

{
\textbf{Étude du signe de $f'(x)$ :}

$\e^{-x} > 0$ pour tout $x \in \mathbf{R}$

Le signe de $f'(x)$ dépend donc de $(4x-7)$ :
\[
4x-7 = 0 \Longleftrightarrow x = \dfrac{7}{4}
\]

\vspace{2mm}

\begin{itemize}
    \item Si $x < \dfrac{7}{4}$ : $4x-7 < 0$ donc \textbf{\textcolor{red}{$f'(x) < 0$}}
    \item Si $x = \dfrac{7}{4}$ : \textbf{\textcolor{red}{$f'(x) = 0$}}
    \item Si $x > \dfrac{7}{4}$ : $4x-7 > 0$ donc \textbf{\textcolor{red}{$f'(x) > 0$}}
\end{itemize}
}
\end{multicols}
\end{tcolorbox}

\vspace{3mm}

\textbf{3. Donner le tableau de variation de $f$.}

\begin{tcolorbox}[colback=red!5!white,colframe=red!75!black]
{
\[
f\left(\dfrac{7}{4}\right) = \e^{-\frac{7}{4}}\left(3-4 \times \dfrac{7}{4}\right) = \e^{-\frac{7}{4}}(3-7) = -4\e^{-\frac{7}{4}}
\]



\begin{center}
\begin{tikzpicture}
    \tkzTabInit[lgt=2,espcl=2.5]
    {$x$ /1, $f'(x)$ /1, $f(x)$ /2}
    {$-\infty$, $\dfrac{7}{4}$, $+\infty$}
    \tkzTabLine{, -, z, +, }
    \tkzTabVar{+/$+\infty$, -/$-4\e^{-\frac{7}{4}}$, +/$0$}
\end{tikzpicture}
\end{center}

}
\end{tcolorbox}

\vspace{3mm}

\textbf{4. Écrire une équation de la tangente à C au point d'abscisse $a=0$}

\begin{tcolorbox}[colback=red!5!white,colframe=red!75!black]
\begin{multicols}{2}
    

{
L'équation de la tangente au point d'abscisse $a$ est : $y = f'(a)(x-a) + f(a)$

\vspace{2mm}

Pour $a = 0$ :
\[
\begin{aligned}
f(0) &= \e^{0}(3-0) = 1 \times 3 = 3 \\
f'(0) &= \e^{0}(0-7) = 1 \times (-7) = -7
\end{aligned}
\]

\vspace{2mm}

\[
\begin{aligned}
y &= f'(0)(x-0) + f(0) \\
y &= -7x + 3
\end{aligned}
\]

\vspace{2mm}

\begin{center}
\textbf{\textcolor{red}{L'équation de la tangente est : $y = -7x + 3$}}
\end{center}
}
\end{multicols}
\end{tcolorbox}


\end{document}
